% --- Archon physics formalization source begin ---
% archon:physics
% archon:covers PhyXMiniProblems/problem_phyx_mini_0837.lean
% archon:source-report reports/phyx_mini/problem_phyx_mini_0837.source.json

\chapter{Physics problem phyx\_mini\_0837}
\label{ch:PhyXMiniProblems_problem_phyx_mini_0837}

\paragraph{Problem source.}
\#\# Physical scenario

The two parallel plates in figure are \textbackslash{}( 2.0\textbackslash{}\textbackslash{},\textbackslash{}\textbackslash{}mathrm\{cm\} \textbackslash{}) apart and the electric field strength between them is \textbackslash{}( 1.0 	imes 10\textasciicircum{}4\textbackslash{}\textbackslash{},\textbackslash{}\textbackslash{}mathrm\{N/C\} \textbackslash{}). An electron is launched at a \textbackslash{}( 45\textasciicircum{}\textbackslash{}circ \textbackslash{}) angle from the positive plate.

\#\# Question

What is the maximum initial speed \textbackslash{}( v\_0 \textbackslash{}) the electron can have without hitting the negative plate?

\#\# Answer choices

- A: A: 1.3 \textbackslash{}times 10\textasciicircum{}\{7\}m/s

- B: B: 5.3 \textbackslash{}times 10\textasciicircum{}\{7\}m/s

- C: C: 1.0 \textbackslash{}times 10\textasciicircum{}\{7\}m/s

- D: D: 1.19 \textbackslash{}times 10\textasciicircum{}\{7\}m/s

\#\# Figure caption (auxiliary; use the image as primary evidence)

The image shows a diagram with two parallel plates. The bottom plate has positive charges indicated by plus signs, while the top plate has a dashed line, possibly indicating neutrality or a different charge.

Between the plates, an arrow labeled \textbackslash{}(\textbackslash{}mathbf\{v\}\_0\textbackslash{}) represents an initial velocity vector aimed upward and to the right, forming a 45° angle with the horizontal. The arrow originates from a small circle at the bottom left, representing a charged particle or object.

There is a vertical double-headed arrow between the plates with text indicating a distance of "2.0 cm" between them. This suggests the particle is moving through a uniform electric field between the plates.

\#\# Dataset metadata

Electrostatics; Physical Model Grounding Reasoning, Multi-Formula Reasoning

\paragraph{Recorded answer/context.}
D: D: 1.19 \textbackslash{}times 10\textasciicircum{}\{7\}m/s

\paragraph{Figure/image path.}
/root/proposal\_for\_physic/science-mango/phyx\_mini\_run/phyx\_data/test\_image/837.png

\paragraph{Formalization target.}
create a compiling Lean file with sorry bodies at `PhyXMiniProblems/problem\_phyx\_mini\_0837.lean`. The Lean declarations must preserve the physical quantities, dimensions or dimensional roles, figure labels, governing-law hypotheses, and final relation expressed by this problem.
Use Mathlib/Physlib names found through LeanExplore where available. If a physics API is missing, introduce faithful local abstractions rather than scalar placeholder aliases.

\begin{theorem}[Physics formalization target]
\label{thm:physics:phyx_mini_0837:target}
\lean{PhyXMiniProblems.ProblemPhyXMini0837.maximum_initial_speed_matches_choiceD}
\uses{def:physics:phyx-mini-0837:phyxminiproblems-problemphyxmini0837-speedinmeterspersecond, def:physics:phyx-mini-0837:phyxminiproblems-problemphyxmini0837-electronparallelplatesetup, def:physics:phyx-mini-0837:phyxminiproblems-problemphyxmini0837-matcheselectronparallelplatescenario, def:physics:phyx-mini-0837:phyxminiproblems-problemphyxmini0837-matchesproblemreadouts, def:physics:phyx-mini-0837:phyxminiproblems-problemphyxmini0837-matchessuppliedparallelplatefigure, def:physics:phyx-mini-0837:phyxminiproblems-problemphyxmini0837-usesstandardelectronreferencedata, def:physics:phyx-mini-0837:phyxminiproblems-problemphyxmini0837-hasphysicalparallelplateparameters, def:physics:phyx-mini-0837:phyxminiproblems-problemphyxmini0837-isuniformupwardparallelplatefield, def:physics:phyx-mini-0837:phyxminiproblems-problemphyxmini0837-satisfieselectronparallelplatemotionlaws, def:physics:phyx-mini-0837:phyxminiproblems-problemphyxmini0837-islimitingnonimpactspeed, def:physics:phyx-mini-0837:phyxminiproblems-problemphyxmini0837-displayedspeedinmeterspersecond, def:physics:phyx-mini-0837:phyxminiproblems-problemphyxmini0837-roundstonearestincrement, def:physics:phyx-mini-0837:phyxminiproblems-problemphyxmini0837-isclosestdisplayedspeed}
The assigned autoformalize agent should translate this physics problem into Lean declarations in the covered file, with theorem and lemma proof bodies written as `by sorry`.
\end{theorem}
\begin{proof}
This is an autoformalization task, not a proof task. Produce statements that can later be proved by the physics prover without weakening the physical model.
\end{proof}

% --- Archon declaration topology begin ---
\section{Declaration topology}
\begin{definition}[acceleration Dimension]
  \label{def:physics:phyx-mini-0837:phyxminiproblems-problemphyxmini0837-accelerationdimension}
  \lean{PhyXMiniProblems.ProblemPhyXMini0837.accelerationDimension}
  The physical dimension of acceleration, L T\textasciicircum{}-2
\end{definition}
\begin{proof}
  This is the declared model definition of acceleration Dimension.
\end{proof}
\begin{definition}[force Dimension]
  \label{def:physics:phyx-mini-0837:phyxminiproblems-problemphyxmini0837-forcedimension}
  \lean{PhyXMiniProblems.ProblemPhyXMini0837.forceDimension}
  \uses{def:physics:phyx-mini-0837:phyxminiproblems-problemphyxmini0837-accelerationdimension}
  The physical dimension of force, M L T\textasciicircum{}-2
\end{definition}
\begin{proof}
  This is the declared model definition of force Dimension.
\end{proof}
\begin{definition}[electric Field Strength Dimension]
  \label{def:physics:phyx-mini-0837:phyxminiproblems-problemphyxmini0837-electricfieldstrengthdimension}
  \lean{PhyXMiniProblems.ProblemPhyXMini0837.electricFieldStrengthDimension}
  \uses{def:physics:phyx-mini-0837:phyxminiproblems-problemphyxmini0837-forcedimension}
  The physical dimension of electric-field strength, M L T\textasciicircum{}-2 C\textasciicircum{}-1
\end{definition}
\begin{proof}
  This is the declared model definition of electric Field Strength Dimension.
\end{proof}
\begin{definition}[Length Quantity]
  \label{def:physics:phyx-mini-0837:phyxminiproblems-problemphyxmini0837-lengthquantity}
  \lean{PhyXMiniProblems.ProblemPhyXMini0837.LengthQuantity}
  A nonnegative, unit-independent physical length
\end{definition}
\begin{proof}
  This is the declared model definition of Length Quantity.
\end{proof}
\begin{definition}[Time Quantity]
  \label{def:physics:phyx-mini-0837:phyxminiproblems-problemphyxmini0837-timequantity}
  \lean{PhyXMiniProblems.ProblemPhyXMini0837.TimeQuantity}
  A nonnegative, unit-independent physical elapsed time
\end{definition}
\begin{proof}
  This is the declared model definition of Time Quantity.
\end{proof}
\begin{definition}[Mass Quantity]
  \label{def:physics:phyx-mini-0837:phyxminiproblems-problemphyxmini0837-massquantity}
  \lean{PhyXMiniProblems.ProblemPhyXMini0837.MassQuantity}
  A nonnegative, unit-independent physical mass
\end{definition}
\begin{proof}
  This is the declared model definition of Mass Quantity.
\end{proof}
\begin{definition}[Signed Charge Quantity]
  \label{def:physics:phyx-mini-0837:phyxminiproblems-problemphyxmini0837-signedchargequantity}
  \lean{PhyXMiniProblems.ProblemPhyXMini0837.SignedChargeQuantity}
  A signed, unit-independent physical electric charge
\end{definition}
\begin{proof}
  This is the declared model definition of Signed Charge Quantity.
\end{proof}
\begin{definition}[Electric Field Strength Quantity]
  \label{def:physics:phyx-mini-0837:phyxminiproblems-problemphyxmini0837-electricfieldstrengthquantity}
  \lean{PhyXMiniProblems.ProblemPhyXMini0837.ElectricFieldStrengthQuantity}
  \uses{def:physics:phyx-mini-0837:phyxminiproblems-problemphyxmini0837-electricfieldstrengthdimension}
  A nonnegative electric-field-strength magnitude
\end{definition}
\begin{proof}
  This is the declared model definition of Electric Field Strength Quantity.
\end{proof}
\begin{definition}[Planar Displacement Quantity]
  \label{def:physics:phyx-mini-0837:phyxminiproblems-problemphyxmini0837-planardisplacementquantity}
  \lean{PhyXMiniProblems.ProblemPhyXMini0837.PlanarDisplacementQuantity}
  A signed planar displacement vector
\end{definition}
\begin{proof}
  This is the declared model definition of Planar Displacement Quantity.
\end{proof}
\begin{definition}[Planar Force Quantity]
  \label{def:physics:phyx-mini-0837:phyxminiproblems-problemphyxmini0837-planarforcequantity}
  \lean{PhyXMiniProblems.ProblemPhyXMini0837.PlanarForceQuantity}
  \uses{def:physics:phyx-mini-0837:phyxminiproblems-problemphyxmini0837-forcedimension}
  A signed planar force vector
\end{definition}
\begin{proof}
  This is the declared model definition of Planar Force Quantity.
\end{proof}
\begin{definition}[Planar Acceleration Quantity]
  \label{def:physics:phyx-mini-0837:phyxminiproblems-problemphyxmini0837-planaraccelerationquantity}
  \lean{PhyXMiniProblems.ProblemPhyXMini0837.PlanarAccelerationQuantity}
  \uses{def:physics:phyx-mini-0837:phyxminiproblems-problemphyxmini0837-accelerationdimension}
  A signed planar acceleration vector
\end{definition}
\begin{proof}
  This is the declared model definition of Planar Acceleration Quantity.
\end{proof}
\begin{definition}[x Axis]
  \label{def:physics:phyx-mini-0837:phyxminiproblems-problemphyxmini0837-xaxis}
  \lean{PhyXMiniProblems.ProblemPhyXMini0837.xAxis}
  Coordinate 0, horizontal and positive to the right in the figure
\end{definition}
\begin{proof}
  This is the declared model definition of x Axis.
\end{proof}
\begin{definition}[y Axis]
  \label{def:physics:phyx-mini-0837:phyxminiproblems-problemphyxmini0837-yaxis}
  \lean{PhyXMiniProblems.ProblemPhyXMini0837.yAxis}
  Coordinate 1, vertical and positive toward the negative plate
\end{definition}
\begin{proof}
  This is the declared model definition of y Axis.
\end{proof}
\begin{definition}[upward Direction]
  \label{def:physics:phyx-mini-0837:phyxminiproblems-problemphyxmini0837-upwarddirection}
  \lean{PhyXMiniProblems.ProblemPhyXMini0837.upwardDirection}
  \uses{def:physics:phyx-mini-0837:phyxminiproblems-problemphyxmini0837-yaxis}
  The unit vector pointing from the positive plate to the negative plate
\end{definition}
\begin{proof}
  This is the declared model definition of upward Direction.
\end{proof}
\begin{definition}[nonnegative SIReadout]
  \label{def:physics:phyx-mini-0837:phyxminiproblems-problemphyxmini0837-nonnegativesireadout}
  \lean{PhyXMiniProblems.ProblemPhyXMini0837.nonnegativeSIReadout}
  Read a nonnegative dimensionful scalar in coherent SI units
\end{definition}
\begin{proof}
  This is the declared model definition of nonnegative SIReadout.
\end{proof}
\begin{definition}[signed SIReadout]
  \label{def:physics:phyx-mini-0837:phyxminiproblems-problemphyxmini0837-signedsireadout}
  \lean{PhyXMiniProblems.ProblemPhyXMini0837.signedSIReadout}
  Read a signed dimensionful scalar in coherent SI units
\end{definition}
\begin{proof}
  This is the declared model definition of signed SIReadout.
\end{proof}
\begin{definition}[planar SIReadout]
  \label{def:physics:phyx-mini-0837:phyxminiproblems-problemphyxmini0837-planarsireadout}
  \lean{PhyXMiniProblems.ProblemPhyXMini0837.planarSIReadout}
  Read a dimensionful planar vector in coherent SI units
\end{definition}
\begin{proof}
  This is the declared model definition of planar SIReadout.
\end{proof}
\begin{definition}[length In Meters]
  \label{def:physics:phyx-mini-0837:phyxminiproblems-problemphyxmini0837-lengthinmeters}
  \lean{PhyXMiniProblems.ProblemPhyXMini0837.lengthInMeters}
  \uses{def:physics:phyx-mini-0837:phyxminiproblems-problemphyxmini0837-lengthquantity, def:physics:phyx-mini-0837:phyxminiproblems-problemphyxmini0837-nonnegativesireadout}
  Metre readout of a physical length
\end{definition}
\begin{proof}
  This is the declared model definition of length In Meters.
\end{proof}
\begin{definition}[time In Seconds]
  \label{def:physics:phyx-mini-0837:phyxminiproblems-problemphyxmini0837-timeinseconds}
  \lean{PhyXMiniProblems.ProblemPhyXMini0837.timeInSeconds}
  \uses{def:physics:phyx-mini-0837:phyxminiproblems-problemphyxmini0837-timequantity, def:physics:phyx-mini-0837:phyxminiproblems-problemphyxmini0837-nonnegativesireadout}
  Second readout of a physical elapsed time
\end{definition}
\begin{proof}
  This is the declared model definition of time In Seconds.
\end{proof}
\begin{definition}[mass In Kilograms]
  \label{def:physics:phyx-mini-0837:phyxminiproblems-problemphyxmini0837-massinkilograms}
  \lean{PhyXMiniProblems.ProblemPhyXMini0837.massInKilograms}
  \uses{def:physics:phyx-mini-0837:phyxminiproblems-problemphyxmini0837-massquantity, def:physics:phyx-mini-0837:phyxminiproblems-problemphyxmini0837-nonnegativesireadout}
  Kilogram readout of a physical mass
\end{definition}
\begin{proof}
  This is the declared model definition of mass In Kilograms.
\end{proof}
\begin{definition}[charge In Coulombs]
  \label{def:physics:phyx-mini-0837:phyxminiproblems-problemphyxmini0837-chargeincoulombs}
  \lean{PhyXMiniProblems.ProblemPhyXMini0837.chargeInCoulombs}
  \uses{def:physics:phyx-mini-0837:phyxminiproblems-problemphyxmini0837-signedchargequantity, def:physics:phyx-mini-0837:phyxminiproblems-problemphyxmini0837-signedsireadout}
  Coulomb readout of a signed physical charge
\end{definition}
\begin{proof}
  This is the declared model definition of charge In Coulombs.
\end{proof}
\begin{definition}[speed In Meters Per Second]
  \label{def:physics:phyx-mini-0837:phyxminiproblems-problemphyxmini0837-speedinmeterspersecond}
  \lean{PhyXMiniProblems.ProblemPhyXMini0837.speedInMetersPerSecond}
  Metre-per-second readout of a physical speed magnitude
\end{definition}
\begin{proof}
  This is the declared model definition of speed In Meters Per Second.
\end{proof}
\begin{definition}[electric Field Strength In Newtons Per Coulomb]
  \label{def:physics:phyx-mini-0837:phyxminiproblems-problemphyxmini0837-electricfieldstrengthinnewtonspercoulomb}
  \lean{PhyXMiniProblems.ProblemPhyXMini0837.electricFieldStrengthInNewtonsPerCoulomb}
  \uses{def:physics:phyx-mini-0837:phyxminiproblems-problemphyxmini0837-electricfieldstrengthquantity, def:physics:phyx-mini-0837:phyxminiproblems-problemphyxmini0837-nonnegativesireadout}
  Newton-per-coulomb readout of an electric-field-strength magnitude
\end{definition}
\begin{proof}
  This is the declared model definition of electric Field Strength In Newtons Per Coulomb.
\end{proof}
\begin{definition}[displacement In Meters]
  \label{def:physics:phyx-mini-0837:phyxminiproblems-problemphyxmini0837-displacementinmeters}
  \lean{PhyXMiniProblems.ProblemPhyXMini0837.displacementInMeters}
  \uses{def:physics:phyx-mini-0837:phyxminiproblems-problemphyxmini0837-planardisplacementquantity, def:physics:phyx-mini-0837:phyxminiproblems-problemphyxmini0837-planarsireadout}
  Cartesian displacement readout in metres
\end{definition}
\begin{proof}
  This is the declared model definition of displacement In Meters.
\end{proof}
\begin{definition}[force In Newtons]
  \label{def:physics:phyx-mini-0837:phyxminiproblems-problemphyxmini0837-forceinnewtons}
  \lean{PhyXMiniProblems.ProblemPhyXMini0837.forceInNewtons}
  \uses{def:physics:phyx-mini-0837:phyxminiproblems-problemphyxmini0837-planarforcequantity, def:physics:phyx-mini-0837:phyxminiproblems-problemphyxmini0837-planarsireadout}
  Cartesian net-force readout in newtons
\end{definition}
\begin{proof}
  This is the declared model definition of force In Newtons.
\end{proof}
\begin{definition}[acceleration In Meters Per Second Squared]
  \label{def:physics:phyx-mini-0837:phyxminiproblems-problemphyxmini0837-accelerationinmeterspersecondsquared}
  \lean{PhyXMiniProblems.ProblemPhyXMini0837.accelerationInMetersPerSecondSquared}
  \uses{def:physics:phyx-mini-0837:phyxminiproblems-problemphyxmini0837-planaraccelerationquantity, def:physics:phyx-mini-0837:phyxminiproblems-problemphyxmini0837-planarsireadout}
  Cartesian acceleration readout in metres per second squared
\end{definition}
\begin{proof}
  This is the declared model definition of acceleration In Meters Per Second Squared.
\end{proof}
\begin{definition}[Particle Species]
  \label{def:physics:phyx-mini-0837:phyxminiproblems-problemphyxmini0837-particlespecies}
  \lean{PhyXMiniProblems.ProblemPhyXMini0837.ParticleSpecies}
  The particle species named in the problem
\end{definition}
\begin{proof}
  This is the declared model definition of Particle Species.
\end{proof}
\begin{definition}[Plate]
  \label{def:physics:phyx-mini-0837:phyxminiproblems-problemphyxmini0837-plate}
  \lean{PhyXMiniProblems.ProblemPhyXMini0837.Plate}
  The two horizontal plates in the supplied raster
\end{definition}
\begin{proof}
  This is the declared model definition of Plate.
\end{proof}
\begin{definition}[Plate Polarity]
  \label{def:physics:phyx-mini-0837:phyxminiproblems-problemphyxmini0837-platepolarity}
  \lean{PhyXMiniProblems.ProblemPhyXMini0837.PlatePolarity}
  Electric polarity attached to each plate
\end{definition}
\begin{proof}
  This is the declared model definition of Plate Polarity.
\end{proof}
\begin{definition}[Launch Arrow Direction]
  \label{def:physics:phyx-mini-0837:phyxminiproblems-problemphyxmini0837-launcharrowdirection}
  \lean{PhyXMiniProblems.ProblemPhyXMini0837.LaunchArrowDirection}
  Qualitative direction of the velocity arrow in the diagram
\end{definition}
\begin{proof}
  This is the declared model definition of Launch Arrow Direction.
\end{proof}
\begin{definition}[Motion Regime]
  \label{def:physics:phyx-mini-0837:phyxminiproblems-problemphyxmini0837-motionregime}
  \lean{PhyXMiniProblems.ProblemPhyXMini0837.MotionRegime}
  Idealized physical regime intended by the textbook problem
\end{definition}
\begin{proof}
  This is the declared model definition of Motion Regime.
\end{proof}
\begin{definition}[Parallel Plate Electron Figure]
  \label{def:physics:phyx-mini-0837:phyxminiproblems-problemphyxmini0837-parallelplateelectronfigure}
  \lean{PhyXMiniProblems.ProblemPhyXMini0837.ParallelPlateElectronFigure}
  \uses{def:physics:phyx-mini-0837:phyxminiproblems-problemphyxmini0837-plate, def:physics:phyx-mini-0837:phyxminiproblems-problemphyxmini0837-platepolarity, def:physics:phyx-mini-0837:phyxminiproblems-problemphyxmini0837-launcharrowdirection}
  Literal and qualitative information transcribed from 837.png. The two real-valued fields are printed figure labels, not solved physical quantities
\end{definition}
\begin{proof}
  This is the declared model definition of Parallel Plate Electron Figure.
\end{proof}
\begin{definition}[Electron Parallel Plate Setup]
  \label{def:physics:phyx-mini-0837:phyxminiproblems-problemphyxmini0837-electronparallelplatesetup}
  \lean{PhyXMiniProblems.ProblemPhyXMini0837.ElectronParallelPlateSetup}
  \uses{def:physics:phyx-mini-0837:phyxminiproblems-problemphyxmini0837-lengthquantity, def:physics:phyx-mini-0837:phyxminiproblems-problemphyxmini0837-timequantity, def:physics:phyx-mini-0837:phyxminiproblems-problemphyxmini0837-massquantity, def:physics:phyx-mini-0837:phyxminiproblems-problemphyxmini0837-signedchargequantity, def:physics:phyx-mini-0837:phyxminiproblems-problemphyxmini0837-electricfieldstrengthquantity, def:physics:phyx-mini-0837:phyxminiproblems-problemphyxmini0837-planardisplacementquantity, def:physics:phyx-mini-0837:phyxminiproblems-problemphyxmini0837-planarforcequantity, def:physics:phyx-mini-0837:phyxminiproblems-problemphyxmini0837-planaraccelerationquantity, def:physics:phyx-mini-0837:phyxminiproblems-problemphyxmini0837-particlespecies, def:physics:phyx-mini-0837:phyxminiproblems-problemphyxmini0837-motionregime, def:physics:phyx-mini-0837:phyxminiproblems-problemphyxmini0837-parallelplateelectronfigure}
  Independent physical quantities and the family of trajectories obtained by varying the launch-speed magnitude. No maximum speed or answer choice is a field of this setup
\end{definition}
\begin{proof}
  This is the declared model definition of Electron Parallel Plate Setup.
\end{proof}
\begin{definition}[Matches Electron Parallel Plate Scenario]
  \label{def:physics:phyx-mini-0837:phyxminiproblems-problemphyxmini0837-matcheselectronparallelplatescenario}
  \lean{PhyXMiniProblems.ProblemPhyXMini0837.MatchesElectronParallelPlateScenario}
  \uses{def:physics:phyx-mini-0837:phyxminiproblems-problemphyxmini0837-electronparallelplatesetup}
  Qualitative physical scenario stated in the problem
\end{definition}
\begin{proof}
  This is the declared model definition of Matches Electron Parallel Plate Scenario.
\end{proof}
\begin{definition}[Matches Problem Readouts]
  \label{def:physics:phyx-mini-0837:phyxminiproblems-problemphyxmini0837-matchesproblemreadouts}
  \lean{PhyXMiniProblems.ProblemPhyXMini0837.MatchesProblemReadouts}
  \uses{def:physics:phyx-mini-0837:phyxminiproblems-problemphyxmini0837-lengthinmeters, def:physics:phyx-mini-0837:phyxminiproblems-problemphyxmini0837-electricfieldstrengthinnewtonspercoulomb, def:physics:phyx-mini-0837:phyxminiproblems-problemphyxmini0837-electronparallelplatesetup}
  The three calibrated physical readouts stated in the prose. This record does not constrain a launch-speed magnitude
\end{definition}
\begin{proof}
  This is the declared model definition of Matches Problem Readouts.
\end{proof}
\begin{definition}[Matches Supplied Parallel Plate Figure]
  \label{def:physics:phyx-mini-0837:phyxminiproblems-problemphyxmini0837-matchessuppliedparallelplatefigure}
  \lean{PhyXMiniProblems.ProblemPhyXMini0837.MatchesSuppliedParallelPlateFigure}
  \uses{def:physics:phyx-mini-0837:phyxminiproblems-problemphyxmini0837-electronparallelplatesetup}
  Facts read directly from the supplied parallel-plate diagram
\end{definition}
\begin{proof}
  This is the declared model definition of Matches Supplied Parallel Plate Figure.
\end{proof}
\begin{definition}[Uses Standard Electron Reference Data]
  \label{def:physics:phyx-mini-0837:phyxminiproblems-problemphyxmini0837-usesstandardelectronreferencedata}
  \lean{PhyXMiniProblems.ProblemPhyXMini0837.UsesStandardElectronReferenceData}
  \uses{def:physics:phyx-mini-0837:phyxminiproblems-problemphyxmini0837-massinkilograms, def:physics:phyx-mini-0837:phyxminiproblems-problemphyxmini0837-chargeincoulombs, def:physics:phyx-mini-0837:phyxminiproblems-problemphyxmini0837-electronparallelplatesetup}
  Standard electron reference data needed by the calculation. The magnitude ratio |q\_e| / m\_e = 1.76 * 10\textasciicircum{}11 C/kg is independent of the requested launch speed and of all displayed choices
\end{definition}
\begin{proof}
  This is the declared model definition of Uses Standard Electron Reference Data.
\end{proof}
\begin{definition}[Has Physical Parallel Plate Parameters]
  \label{def:physics:phyx-mini-0837:phyxminiproblems-problemphyxmini0837-hasphysicalparallelplateparameters}
  \lean{PhyXMiniProblems.ProblemPhyXMini0837.HasPhysicalParallelPlateParameters}
  \uses{def:physics:phyx-mini-0837:phyxminiproblems-problemphyxmini0837-lengthinmeters, def:physics:phyx-mini-0837:phyxminiproblems-problemphyxmini0837-massinkilograms, def:physics:phyx-mini-0837:phyxminiproblems-problemphyxmini0837-chargeincoulombs, def:physics:phyx-mini-0837:phyxminiproblems-problemphyxmini0837-electricfieldstrengthinnewtonspercoulomb, def:physics:phyx-mini-0837:phyxminiproblems-problemphyxmini0837-electronparallelplatesetup}
  Positivity and angle-range conditions selecting the physical branch
\end{definition}
\begin{proof}
  This is the declared model definition of Has Physical Parallel Plate Parameters.
\end{proof}
\begin{definition}[Is Uniform Upward Parallel Plate Field]
  \label{def:physics:phyx-mini-0837:phyxminiproblems-problemphyxmini0837-isuniformupwardparallelplatefield}
  \lean{PhyXMiniProblems.ProblemPhyXMini0837.IsUniformUpwardParallelPlateField}
  \uses{def:physics:phyx-mini-0837:phyxminiproblems-problemphyxmini0837-upwarddirection, def:physics:phyx-mini-0837:phyxminiproblems-problemphyxmini0837-electricfieldstrengthinnewtonspercoulomb, def:physics:phyx-mini-0837:phyxminiproblems-problemphyxmini0837-electronparallelplatesetup}
  The Physlib electric field is spatially and temporally uniform and points from the positive lower plate toward the negative upper plate. Its vector entries are interpreted as coherent-SI field-strength components
\end{definition}
\begin{proof}
  This is the declared model definition of Is Uniform Upward Parallel Plate Field.
\end{proof}
\begin{definition}[Satisfies Electron Parallel Plate Motion Laws]
  \label{def:physics:phyx-mini-0837:phyxminiproblems-problemphyxmini0837-satisfieselectronparallelplatemotionlaws}
  \lean{PhyXMiniProblems.ProblemPhyXMini0837.SatisfiesElectronParallelPlateMotionLaws}
  \uses{def:physics:phyx-mini-0837:phyxminiproblems-problemphyxmini0837-xaxis, def:physics:phyx-mini-0837:phyxminiproblems-problemphyxmini0837-yaxis, def:physics:phyx-mini-0837:phyxminiproblems-problemphyxmini0837-timeinseconds, def:physics:phyx-mini-0837:phyxminiproblems-problemphyxmini0837-massinkilograms, def:physics:phyx-mini-0837:phyxminiproblems-problemphyxmini0837-chargeincoulombs, def:physics:phyx-mini-0837:phyxminiproblems-problemphyxmini0837-speedinmeterspersecond, def:physics:phyx-mini-0837:phyxminiproblems-problemphyxmini0837-electricfieldstrengthinnewtonspercoulomb, def:physics:phyx-mini-0837:phyxminiproblems-problemphyxmini0837-displacementinmeters, def:physics:phyx-mini-0837:phyxminiproblems-problemphyxmini0837-forceinnewtons, def:physics:phyx-mini-0837:phyxminiproblems-problemphyxmini0837-accelerationinmeterspersecondsquared, def:physics:phyx-mini-0837:phyxminiproblems-problemphyxmini0837-electronparallelplatesetup}
  Governing force and kinematic laws: the horizontal electric-force component vanishes; the vertical component is q E in the upward-positive convention; Newton's second law holds componentwise; and every trial speed produces the usual constant-acceleration trajectory from the launch point on the lower plate. These equations apply uniformly to arbitrary trial speeds and contain no limiting speed or displayed answer value
\end{definition}
\begin{proof}
  This is the declared model definition of Satisfies Electron Parallel Plate Motion Laws.
\end{proof}
\begin{definition}[Avoids Negative Plate]
  \label{def:physics:phyx-mini-0837:phyxminiproblems-problemphyxmini0837-avoidsnegativeplate}
  \lean{PhyXMiniProblems.ProblemPhyXMini0837.AvoidsNegativePlate}
  \uses{def:physics:phyx-mini-0837:phyxminiproblems-problemphyxmini0837-yaxis, def:physics:phyx-mini-0837:phyxminiproblems-problemphyxmini0837-lengthinmeters, def:physics:phyx-mini-0837:phyxminiproblems-problemphyxmini0837-displacementinmeters, def:physics:phyx-mini-0837:phyxminiproblems-problemphyxmini0837-electronparallelplatesetup}
  A trial launch avoids the negative plate when its vertical coordinate remains strictly below the plate separation at every nonnegative physical elapsed time
\end{definition}
\begin{proof}
  This is the declared model definition of Avoids Negative Plate.
\end{proof}
\begin{definition}[Is Limiting Nonimpact Speed]
  \label{def:physics:phyx-mini-0837:phyxminiproblems-problemphyxmini0837-islimitingnonimpactspeed}
  \lean{PhyXMiniProblems.ProblemPhyXMini0837.IsLimitingNonimpactSpeed}
  \uses{def:physics:phyx-mini-0837:phyxminiproblems-problemphyxmini0837-speedinmeterspersecond, def:physics:phyx-mini-0837:phyxminiproblems-problemphyxmini0837-electronparallelplatesetup, def:physics:phyx-mini-0837:phyxminiproblems-problemphyxmini0837-avoidsnegativeplate}
  speedLimit is the supremal non-impact speed: it bounds every avoiding speed, and every strictly slower physical speed avoids the upper plate. Membership of the threshold itself is intentionally not required
\end{definition}
\begin{proof}
  This is the declared model definition of Is Limiting Nonimpact Speed.
\end{proof}
\begin{definition}[Answer Choice]
  \label{def:physics:phyx-mini-0837:phyxminiproblems-problemphyxmini0837-answerchoice}
  \lean{PhyXMiniProblems.ProblemPhyXMini0837.AnswerChoice}
  Labels attached to the four displayed speed choices
\end{definition}
\begin{proof}
  This is the declared model definition of Answer Choice.
\end{proof}
\begin{definition}[displayed Speed In Meters Per Second]
  \label{def:physics:phyx-mini-0837:phyxminiproblems-problemphyxmini0837-displayedspeedinmeterspersecond}
  \lean{PhyXMiniProblems.ProblemPhyXMini0837.displayedSpeedInMetersPerSecond}
  \uses{def:physics:phyx-mini-0837:phyxminiproblems-problemphyxmini0837-answerchoice}
  Metre-per-second value printed beside each answer label
\end{definition}
\begin{proof}
  This is the declared model definition of displayed Speed In Meters Per Second.
\end{proof}
\begin{definition}[Rounds To Nearest Increment]
  \label{def:physics:phyx-mini-0837:phyxminiproblems-problemphyxmini0837-roundstonearestincrement}
  \lean{PhyXMiniProblems.ProblemPhyXMini0837.RoundsToNearestIncrement}
  A general nearest-rounding relation at a positive displayed increment
\end{definition}
\begin{proof}
  This is the declared model definition of Rounds To Nearest Increment.
\end{proof}
\begin{definition}[Is Closest Displayed Speed]
  \label{def:physics:phyx-mini-0837:phyxminiproblems-problemphyxmini0837-isclosestdisplayedspeed}
  \lean{PhyXMiniProblems.ProblemPhyXMini0837.IsClosestDisplayedSpeed}
  \uses{def:physics:phyx-mini-0837:phyxminiproblems-problemphyxmini0837-speedinmeterspersecond, def:physics:phyx-mini-0837:phyxminiproblems-problemphyxmini0837-answerchoice, def:physics:phyx-mini-0837:phyxminiproblems-problemphyxmini0837-displayedspeedinmeterspersecond}
  A choice is closest when no displayed alternative is nearer
\end{definition}
\begin{proof}
  This is the declared model definition of Is Closest Displayed Speed.
\end{proof}
\begin{lemma}[limiting nonimpact speed squared formula]
  \label{lem:physics:phyx-mini-0837:phyxminiproblems-problemphyxmini0837-limiting-nonimpact-speed-squared-formula}
  \lean{PhyXMiniProblems.ProblemPhyXMini0837.limiting_nonimpact_speed_squared_formula}
  \uses{def:physics:phyx-mini-0837:phyxminiproblems-problemphyxmini0837-lengthinmeters, def:physics:phyx-mini-0837:phyxminiproblems-problemphyxmini0837-massinkilograms, def:physics:phyx-mini-0837:phyxminiproblems-problemphyxmini0837-chargeincoulombs, def:physics:phyx-mini-0837:phyxminiproblems-problemphyxmini0837-speedinmeterspersecond, def:physics:phyx-mini-0837:phyxminiproblems-problemphyxmini0837-electricfieldstrengthinnewtonspercoulomb, def:physics:phyx-mini-0837:phyxminiproblems-problemphyxmini0837-electronparallelplatesetup, def:physics:phyx-mini-0837:phyxminiproblems-problemphyxmini0837-usesstandardelectronreferencedata, def:physics:phyx-mini-0837:phyxminiproblems-problemphyxmini0837-hasphysicalparallelplateparameters, def:physics:phyx-mini-0837:phyxminiproblems-problemphyxmini0837-isuniformupwardparallelplatefield, def:physics:phyx-mini-0837:phyxminiproblems-problemphyxmini0837-satisfieselectronparallelplatemotionlaws, def:physics:phyx-mini-0837:phyxminiproblems-problemphyxmini0837-islimitingnonimpactspeed}
  For an arbitrary physical instance, the vertical trajectory and electric force laws give v\_limit\textasciicircum{}2 = 2 (|q\_e| / m\_e) E d / sin(theta)\textasciicircum{}2. This symbolic relation is derived from the governing laws and limiting property; it does not occur in a premise structure
\end{lemma}
\begin{proof}
  The conclusion follows from the governing relations and figure readouts named in the statement of limiting nonimpact speed squared formula.
\end{proof}
% --- Archon declaration topology end ---
% --- Archon physics formalization source end ---
